\documentclass[11pt,german,hideothersubsections]{beamer}\usepackage[]{graphicx}\usepackage[]{color}
%% maxwidth is the original width if it is less than linewidth
%% otherwise use linewidth (to make sure the graphics do not exceed the margin)
\makeatletter
\def\maxwidth{ %
  \ifdim\Gin@nat@width>\linewidth
    \linewidth
  \else
    \Gin@nat@width
  \fi
}
\makeatother

\definecolor{fgcolor}{rgb}{0.345, 0.345, 0.345}
\newcommand{\hlnum}[1]{\textcolor[rgb]{0.686,0.059,0.569}{#1}}%
\newcommand{\hlstr}[1]{\textcolor[rgb]{0.192,0.494,0.8}{#1}}%
\newcommand{\hlcom}[1]{\textcolor[rgb]{0.678,0.584,0.686}{\textit{#1}}}%
\newcommand{\hlopt}[1]{\textcolor[rgb]{0,0,0}{#1}}%
\newcommand{\hlstd}[1]{\textcolor[rgb]{0.345,0.345,0.345}{#1}}%
\newcommand{\hlkwa}[1]{\textcolor[rgb]{0.161,0.373,0.58}{\textbf{#1}}}%
\newcommand{\hlkwb}[1]{\textcolor[rgb]{0.69,0.353,0.396}{#1}}%
\newcommand{\hlkwc}[1]{\textcolor[rgb]{0.333,0.667,0.333}{#1}}%
\newcommand{\hlkwd}[1]{\textcolor[rgb]{0.737,0.353,0.396}{\textbf{#1}}}%

\usepackage{framed}
\makeatletter
\newenvironment{kframe}{%
 \def\at@end@of@kframe{}%
 \ifinner\ifhmode%
  \def\at@end@of@kframe{\end{minipage}}%
  \begin{minipage}{\columnwidth}%
 \fi\fi%
 \def\FrameCommand##1{\hskip\@totalleftmargin \hskip-\fboxsep
 \colorbox{shadecolor}{##1}\hskip-\fboxsep
     % There is no \\@totalrightmargin, so:
     \hskip-\linewidth \hskip-\@totalleftmargin \hskip\columnwidth}%
 \MakeFramed {\advance\hsize-\width
   \@totalleftmargin\z@ \linewidth\hsize
   \@setminipage}}%
 {\par\unskip\endMakeFramed%
 \at@end@of@kframe}
\makeatother

\definecolor{shadecolor}{rgb}{.97, .97, .97}
\definecolor{messagecolor}{rgb}{0, 0, 0}
\definecolor{warningcolor}{rgb}{1, 0, 1}
\definecolor{errorcolor}{rgb}{1, 0, 0}
\newenvironment{knitrout}{}{} % an empty environment to be redefined in TeX

\usepackage{alltt}

\usepackage{hyperref}
\usepackage{amsmath,nicefrac,booktabs,mathabx}
\usepackage{natbib}
\usepackage{url}
\usepackage{textpos}
\usepackage{listings}
\definecolor{Rblau}{rgb}{.3,.6,.9}

\lstset{language=R,
basicstyle=\ttfamily\footnotesize,
keywordstyle=\color{blue}\bfseries,
identifierstyle=\color{Rblau},
commentstyle=\color{gray},
stringstyle=\color{green}\ttfamily,
showstringspaces=false,
frame=tb}



\bibpunct{(}{)}{;}{a}{,}{,}
\usepackage[english]{babel}
\usepackage[latin1]{inputenc}
\usepackage{helvet}
\usepackage{graphicx}
\usepackage{color}
\usepackage{multirow,dcolumn}
\usepackage{ragged2e}
\usepackage{xcolor}
\usepackage{colortbl}
\usepackage{tikz}
\usetikzlibrary{calc}
\usepackage{booktabs}
\colorlet{tablesubheadcolor}{gray!25}
\colorlet{tableheadcolor}{gray!40}
\colorlet{tablerowcolor}{gray!15.0}
\usetheme[english]{Gesis}
\setbeamertemplate{navigation symbols}{}
\setbeamertemplate{footline}[frame number]%{\hspace*{.2cm}\insertframenumber}
\setbeamerfont{caption}{size=\footnotesize}
\usefonttheme[onlylarge]{structuresmallcapsserif} % alte Schrift

\newcommand{\R}[1]{{\tt \color{blue}  #1}}
  \newtheorem{thm}{Theorem}
  \newtheorem{rem}{Bemerkung}
  \newtheorem{lem}{Lemma}
  
  \definecolor{hellgrau}{rgb}   {0.109375,  0.40625,   0.51953125}
  \definecolor{dunkelgrau}{rgb} {0.009375,  0.30625,   0.41953125}
  \definecolor{dunkelgrau2}{rgb}{0.009375,  0.20625,   0.31953125}
  \definecolor{hellbraun}{rgb}  {0.9140625, 0.8984375, 0.8046875}
  \definecolor{hellbraun2}{rgb} {.95,       0.9,       0.8}
  \definecolor{alertred}{rgb}   {0.8515625, 0.3828125, 0.08984375}
  \definecolor{orange}{rgb}{1,0.5,0}
  
  
  \setbeamercolor{firstsecslide}{fg=white,bg=dunkelgrau}
  \setbeamertemplate{blocks}[rounded][shadow=true]
  
  \newcolumntype{d}[0]{D{,}{.}{6}}
  
  \newenvironment{itemizeol}{\begin{itemize}[<+->]}{\end{itemize}}
  \newenvironment{descriptionol}{\begin{description}[<+->]}{\end{description}}
  
  \newcolumntype{V}[1]{ {\RaggedRight\hspace{0pt}}p{#1}}

\newcommand{\emphred}[1]{\textcolor{alertred}{#1}}
\newcommand{\emphcol}[1]{\textcolor{dunkelgrau}{\slshape #1}}
 
 \setcounter{tocdepth}{1}
 \setbeamercolor*{section in toc}{fg=hellgrau}
\setbeamertemplate{bibliography item}[default]
 \makeatother
\addtobeamertemplate{frametitle}{}{%
 \begin{textblock*}{100mm}(.91\textwidth,-1cm)
 \includegraphics[height=1cm,width=2cm]{../common/pics/GESIS_Logo_kompakt_en.jpg}
 \end{textblock*}}
 \title[Day 1]{Sampling, Weighting and Estimation\\ \Large{Exercise 1} }
 %\subtitle{Umgang am Beispiel von Telefonstichproben}
 
 \author[M. Sand]{Stefan Zins, Matthias Sand\\ and Jan-Philipp Kolb\\ \vspace{.5cm} \footnotesize{GESIS - Leibniz Institute\\ for the Social Sciences}}
 %\institute{\includegraphics[width=4.5cm]{GESIS_Logo_informell}}
%  \date[]{\color{dunkelgrau}\footnotesize%
% \begin{minipage}{8cm}%
% \begin{center}%
%  \scriptsize{
% \textbf{GESIS Summer School}\\ \tiny{Cologne, Germany}%
%  }\\
% \vspace{0.25cm}
%  \textbf{August 24th, 2015}%
%  
%  \end{center}%
%  \end{minipage}}%
\IfFileExists{upquote.sty}{\usepackage{upquote}}{}
\begin{document}




\maketitle



%%%%%%%%%%%%%%%%%%%%%%%%%%%%%%%%%%%%%%%%%%%%%%%%%%
\begin{frame}[fragile]{Exercise 1}

\begin{enumerate}

\item Download the ESS dataset for Sweden (Sampling Data and Country File) of the 5th round
\item[] \alert{\url{http://www.europeansocialsurvey.org/data/country.html?c=sweden}}
\item Setup your workspace and load the R-packages \R{foreign} and \R{survey}
\item Load the ESS dataset and the country file
\item Merge both data frames by their ID-variable, using the \R{merge()}-command

\end{enumerate}

\end{frame}
%%%%%%%%%%%%%%%%%%%%%%%%%%%%%%%%%%%%%%%%%%%%%%%%%%%%%
\begin{frame}[fragile]{Exercise 1}

\begin{enumerate}

\setcounter{enumi}{4}
\item Determine the sampling strategy (Inspect the variables \alert{SAMPPOIN},\alert{STRATFY} and \alert{PROB})
\item Add the variable \R{N}$=\R{dweight}*\R{pweight}*10000*n$ to your data frame that equals the total population size for the target population of the ESS Round 5 of Sweden
\item Convert your data frame to a survey object (\R{survey}-Package)
\item Estimate the total and mean of the variable \R{tvtot}

\end{enumerate}

\end{frame}
%%%%%%%%%%%%%%%%%%%%%%%%%%%%%%%%%%%%%%%%%%%%%%%%%%%%%%
\begin{frame}[fragile]{The \R{survey}-Package}
\footnotesize{
\begin{itemize}
\item The survey package provides a large range of applications for complex survey samples
\item Typically, the first step is to define a survey object with the \R{svydesign()} command
\end{itemize}
\begin{center}
\textbf{Simple survey object (Simple Random Sample)}
\end{center}




\begin{knitrout}\footnotesize
\definecolor{shadecolor}{rgb}{0.969, 0.969, 0.969}\color{fgcolor}\begin{kframe}
\begin{alltt}
\hlkwd{data}\hlstd{(api)}

\hlstd{surv.obj} \hlkwb{<-} \hlkwd{svydesign}\hlstd{(}\hlkwc{id}\hlstd{=}\hlopt{~}\hlnum{1}\hlstd{,}\hlkwc{fpc} \hlstd{=} \hlopt{~}\hlstd{fpc,} \hlkwc{data} \hlstd{= apisrs)}
\end{alltt}
\end{kframe}
\end{knitrout}
}
\end{frame}
%%%%%%%%%%%%%%%%%%%%%%%%%%%%%%%%%%%%%%%%%%%%%%%%%%%%%%%%%
\begin{frame}[fragile]{The survey package}
%\frametitle{\vspace{-.05cm}\begin{center}\footnotesize{The survey package}\end{center}}
\footnotesize{
\begin{knitrout}\footnotesize
\definecolor{shadecolor}{rgb}{0.969, 0.969, 0.969}\color{fgcolor}\begin{kframe}
\begin{alltt}
\hlstd{surv.obj} \hlkwb{<-} \hlkwd{svydesign}\hlstd{(}\hlkwc{id}\hlstd{=}\hlopt{~}\hlnum{1}\hlstd{,}\hlkwc{fpc} \hlstd{=} \hlopt{~}\hlstd{fpc,} \hlkwc{data} \hlstd{= apisrs)}
\end{alltt}
\end{kframe}
\end{knitrout}
}
\begin{itemize}
\item \R{id} specifies the identifier of PSU and SSU;\R{id$=$ \textasciitilde0} or \R{id$=$\textasciitilde1} stipulates a single stage sampling
\item For multi-stage samples the \R{id} argument should always specify a formula with the cluster-identifier at each stage
\item \R{fpc} should be used for the finite population correction
\item[$\Rightarrow$] Either as the total population size of each stratum or as a fraction of the total population that has been sampled 
\item \R{data} reflects the data set for which the design object should be defined
\end{itemize}
\end{frame}
%%%%%%%%%%%%%%%%%%%%%%%%%%%%%%%%%%%%%%%%%%%%%%%%%%%%%%%%%
%%%%%%%%%%%%%%%%%%%%%%%%%%%%%%%%%%%%%%%%%%%%%%%%%%%%%%%%%
\begin{frame}[fragile]{The survey package}
%\frametitle{\vspace{-.05cm}\begin{center}\footnotesize{The survey package}\end{center}}
\footnotesize{
\begin{center}
\textbf{Important commands}
\end{center}
\begin{tabular}{ll}
\R{svytotal} & returns the estimated total of a variable \\ & and its standard error ($+ deff$)\\
\R{svymean} & returns the estimated mean of a variable\\ & and its standard error ($+ deff$)\\
\R{svyquantile} & Computes quantiles for data from complex surveys\\
\R{svyvar} & Computes variances  for data from complex surveys\\
\R{weights} & Returns the (design) weights of a survey object\\
\R{calibrate} & Calibration of a data set (uses the GREG-Estimator)\\
... & ...\\
\end{tabular}

\begin{knitrout}\footnotesize
\definecolor{shadecolor}{rgb}{0.969, 0.969, 0.969}\color{fgcolor}\begin{kframe}
\begin{alltt}
\hlkwd{svytotal}\hlstd{(}\hlopt{~}\hlstd{api00,surv.obj)}
\end{alltt}
\begin{verbatim}
##         total    SE
## api00 4066888 57293
\end{verbatim}
\end{kframe}
\end{knitrout}
}
\end{frame}




\end{document}
